
\documentclass{article}
\usepackage{colortbl}
\usepackage{makecell}
\usepackage{multirow}
\usepackage{supertabular}

\begin{document}

\newcounter{utterance}

\centering \large Interaction Transcript for game `ifeval', experiment `ifeval', episode 240 with llama{-}3.1{-}8b{-}instruct.
\vspace{24pt}

{ \footnotesize  \setcounter{utterance}{1}
\setlength{\tabcolsep}{0pt}
\begin{supertabular}{c@{$\;$}|p{.15\linewidth}@{}p{.15\linewidth}p{.15\linewidth}p{.15\linewidth}p{.15\linewidth}p{.15\linewidth}}
    \# & \multicolumn{2}{c}{Player} && \multicolumn{2}{c}{Game Master} \\
    \hline

    \theutterance \stepcounter{utterance}  
    & & & \multicolumn{4}{p{0.6\linewidth}}{
        \cellcolor[rgb]{0.9,0.9,0.9}{
            \makecell[{{p{\linewidth}}}]{
                \texttt{\tiny{[P1$\langle$GM]}}
                \texttt{Write a casual blog post about how the outer solar system is different from the inner solar system, and what that means for the possibility of life. Wrap your entire response with double quotation marks. Your response should contain 17 or more sentences.} \\
            }
        }
    }
     \\ \\

    \theutterance \stepcounter{utterance}  
    & \multicolumn{4}{p{0.6\linewidth}}{
        \cellcolor[rgb]{0.561,0.737,0.561}{
            \makecell[{{p{\linewidth}}}]{
                \texttt{\tiny{[P1$\rangle$GM]}}
                \texttt{"Exploring the vastness of our solar system, we often find ourselves comparing the inner and outer regions. While the inner solar system is home to the rocky planets like Earth, Mercury, and Mars, the outer solar system is a realm of gas giants, icy worlds, and mysterious dwarf planets. The most striking difference between the two is the temperature and composition of the planets.} \\
\\ 
\texttt{In the inner solar system, the temperatures are relatively warm, with the average temperature on Mercury reaching up to 800 degrees Fahrenheit during the day. This is due to the close proximity to the sun, which provides a significant amount of solar radiation. In contrast, the outer solar system is much colder, with temperatures ranging from {-}200 to {-}400 degrees Fahrenheit on some of the gas giants like Neptune and Uranus. This extreme cold makes it difficult for liquid water to exist, a crucial ingredient for life as we know it.} \\
\\ 
\texttt{Another key difference is the composition of the planets. The inner solar system is dominated by rocky planets with iron{-}rich cores, while the outer solar system is home to gas giants like Jupiter and Saturn, which are primarily composed of hydrogen and helium. The gas giants also have numerous moons, some of which are thought to have conditions suitable for life. For example, Europa, a moon of Jupiter, has a thick icy crust covering a liquid water ocean, making it a prime target in the search for extraterrestrial life.} \\
\\ 
\texttt{The presence of liquid water is a crucial factor in the possibility of life. Water is essential for life as we know it, and its presence on a planet or moon is often seen as a sign of potential habitability. In the outer solar system, the gas giants and their moons offer a unique opportunity to search for life in the form of liquid water. For instance, Enceladus, a moon of Saturn, has geysers of water vapor and organic compounds, suggesting that it may have a liquid water ocean beneath its surface.} \\
\\ 
\texttt{The outer solar system also offers a glimpse into the early days of our solar system. The gas giants and their moons are thought to have formed from a disk of gas and dust that surrounded the sun after its formation. This disk, known as the solar nebula, is believed to have been the source of the building blocks of the planets. By studying the outer solar system, scientists can gain insights into the early history of our solar system and the processes that shaped the planets.} \\
\\ 
\texttt{In conclusion, the outer solar system is a vastly different region from the inner solar system, with its own unique characteristics and opportunities for life. While the inner solar system is home to rocky planets with iron{-}rich cores, the outer solar system is dominated by gas giants and their moons, which offer a chance to search for life in the form of liquid water. As we continue to explore the outer solar system, we may uncover new insights into the possibility of life beyond Earth and the early history of our solar system."} \\
            }
        }
    }
    & & \\ \\

    \theutterance \stepcounter{utterance}  
    & & & \multicolumn{2}{p{0.3\linewidth}}{
        \cellcolor[rgb]{0.9,0.9,0.9}{
            \makecell[{{p{\linewidth}}}]{
                \texttt{\tiny{[GM$|$GM]}}
                {'startend:quotation': {}, 'length_constraints:number_sentences': {'relation': 'at least', 'num_sentences': 17}}
            }
        }
    }
    & & \\ \\

    \theutterance \stepcounter{utterance}  
    & & & \multicolumn{2}{p{0.3\linewidth}}{
        \cellcolor[rgb]{0.9,0.9,0.9}{
            \makecell[{{p{\linewidth}}}]{
                \texttt{\tiny{[GM$|$GM]}}
                \texttt{game\_result = WIN} \\
            }
        }
    }
    & & \\ \\

\end{supertabular}
}

\end{document}
